\documentclass[a4paper,10pt]{amsart}

\pdfinfoomitdate=1
\pdftrailerid{}
\pdfsuppressptexinfo=15

\usepackage[T1]{fontenc}
\usepackage{lmodern}
\usepackage[margin=27mm]{geometry}
\usepackage{amsmath,amssymb,mathtools,booktabs,microtype,xcolor}
\usepackage[numbers,sort&compress]{natbib}
\usepackage[colorlinks=true,linkcolor=blue!55!black,citecolor=blue!55!black,urlcolor=blue!55!black]{hyperref}
\hypersetup{pdftitle={},pdfauthor={},pdfsubject={},pdfkeywords={}}

\newtheorem{theorem}{Theorem}
\newtheorem{proposition}[theorem]{Proposition}
\newtheorem{lemma}[theorem]{Lemma}
\theoremstyle{remark}
\newtheorem{remark}[theorem]{Remark}

\newcommand{\Comp}{\mathop{\mathrm{Comp}}}
\newcommand{\Fix}{\mathop{\mathrm{Fix}}}

\title[Adjacent-run consolidation]{Adjacent-Run Consolidation of Integer Compositions:\
A Sharp Logarithmic Clock and Divisor-Path Fibres}
\author{Anonymous}
\date{}

\begin{document}

\begin{abstract}
Fix a positive integer $n$.  In a composition of $n$, replace every maximal
run of $r$ equal parts $s$ simultaneously by the single part $rs$, and
iterate.  This weight-preserving map is absorbed at the adjacent-unequal
(Carlitz) compositions.  We prove that its largest absorption time is exactly
$\lfloor\log_2 n\rfloor$ for every $n$, and give an explicit equality witness
for every carrier size.  We also determine every one-step fibre, refined by
source length: for a target $\beta=(b_1,\ldots,b_k)$, its fibre polynomial is
the sum of $u^{\sum_i b_i/s_i}$ over divisor choices $s_i\mid b_i$ satisfying
$s_i\ne s_{i+1}$.  Thus the inverse is a weighted path problem on adjacent
divisor sets.  An exhaustive exact-arithmetic falsifier checks all $262{,}143$
compositions of total at most $18$ and performs $2{,}690{,}869$ assertions.
Classical Carlitz enumeration and static run statistics are treated as
background, not as contributions.  The note is anonymous and remains under
\texttt{HOLD\_EXTERNAL}.
\end{abstract}

\maketitle

\section{The finite map and the theorem}

Let $\Comp(n)$ be the set of ordered tuples of positive integers summing to
$n$.  Every $\alpha\in\Comp(n)$ has a unique maximal-run factorisation
\[
  \alpha=s_1^{r_1}s_2^{r_2}\cdots s_k^{r_k},
  \qquad s_i\ne s_{i+1},
\]
where $s^r$ means $r$ adjacent copies of $s$.  Define
\begin{equation}\label{eq:map}
  A_n(\alpha)=(r_1s_1,r_2s_2,\ldots,r_ks_k).
\end{equation}
The update is simultaneous.  For example,
\[
 (1,1,2,4)\longmapsto(2,2,4)\longmapsto(4,4)
 \longmapsto(8).
\]
Write $\tau(\alpha)$ for the first time at which the orbit becomes fixed.

For a target $\beta=(b_1,\ldots,b_k)\in\Comp(n)$, introduce the polynomial
\begin{equation}\label{eq:fibrepoly}
 \Phi_\beta(u)=
 \sum_{\substack{s_i\in\mathbb Z_{>0},\ s_i\mid b_i\ (1\le i\le k)\\
                  s_i\ne s_{i+1}\ (1\le i<k)}}
 u^{\sum_{i=1}^k b_i/s_i}.
\end{equation}
The empty sum is zero.

\begin{theorem}[clock and complete one-step inverse]\label{thm:main}
For every $n\ge1$, the following statements hold.
\begin{enumerate}
\item The fixed points of $A_n$ are exactly the compositions with adjacent
parts unequal.  Every orbit is absorbed at one of them.
\item The sharp absorption clock is
\begin{equation}\label{eq:clock}
 \max_{\alpha\in\Comp(n)}\tau(\alpha)=\lfloor\log_2 n\rfloor.
\end{equation}
An explicit equality witness exists for every $n$.
\item For every target $\beta\in\Comp(n)$ and every $\ell\ge1$,
\begin{equation}\label{eq:coefficient}
 [u^\ell]\Phi_\beta(u)
 =\#\{\alpha\in A_n^{-1}(\beta):\text{$\alpha$ has $\ell$ parts}\}.
\end{equation}
In particular, $|A_n^{-1}(\beta)|=\Phi_\beta(1)$, and $\beta$ lies in the
image precisely when $\Phi_\beta$ is nonzero.
\end{enumerate}
\end{theorem}

The fixed class in part (1) is the classical class of Carlitz compositions
\cite{KnopfmacherProdinger1998}.  Runs of equal parts, their bounded lengths,
and their asymptotics have been studied in considerably greater generality
\cite{Wilf2011,Gafni2015,BenderCanfield2009}.  None of those static results is
claimed here.  The object of the note is the iterated, weight-preserving map
\eqref{eq:map}, its sharp all-size clock, and its target-resolved inverse.

\section{Termination and the sharp clock}

The total weight is visibly invariant under $A_n$.  If a state is not fixed,
at least one run has length at least two, and replacing all runs decreases the
number of parts strictly.  Thus every orbit terminates.  A sharper argument
uses a backwards chain of newly created equalities.

\begin{lemma}[doubling ancestry]\label{lem:doubling}
If an orbit performs $t$ nonfixed updates, then its total weight is at least
$2^t$.
\end{lemma}

\begin{proof}
Put $\alpha^{(j)}=A_n^j(\alpha)$ for $0\le j\le t$.  There is nothing to
prove when $t=0$, so suppose $t\ge1$.  We first isolate the one-step ancestry
claim used below.  If a nontrivial maximal run in $\alpha^{(j)}$, with
$j\ge1$, is collapsed at the next update, choose two adjacent parts of that
run.  They are adjacent outputs of two consecutive maximal runs of
$\alpha^{(j-1)}$.  Those two preceding runs cannot both be singletons: if
they were, their unchanged output values would be equal and the two runs
would instead form one maximal run.  Thus at least one chosen input part was
itself produced by a nontrivial collapse in the preceding update.

Choose a nontrivial run in $\alpha^{(t-1)}$ and call its output in
$\alpha^{(t)}$ the part $w_t$.  When $t\ge2$, apply the ancestry claim to
choose among its equal inputs a part $w_{t-1}$ that was produced
nontrivially from $\alpha^{(t-2)}$; repeat down to a nontrivially produced
part $w_1$ of $\alpha^{(1)}$.  For $2\le j\le t$, the run producing $w_j$
has at least two inputs, all equal to the selected input $w_{j-1}$, and hence
$w_j\ge2w_{j-1}$.  The first selected output satisfies $w_1\ge2$ because it
combines at least two positive original parts.  Therefore
\[
                     2^t\le w_t\le n,
\]
the last inequality holding because $w_t$ is one part of a composition of
total weight $n$.
\end{proof}

Lemma~\ref{lem:doubling} gives
$\tau(\alpha)\le\lfloor\log_2 n\rfloor$.  It remains to attain the bound for
every $n$, including values between powers of two.

For $t\ge1$ put
\begin{equation}\label{eq:cascade}
 C_t=(1,1,2,4,\ldots,2^{t-1}),
\end{equation}
which has weight $2^t$.  After $j<t$ updates, its active prefix is
$(2^j,2^j)$, followed by $2^{j+1},\ldots,2^{t-1}$; hence it takes exactly
$t$ updates to reach $(2^t)$.

Now let $t=\lfloor\log_2n\rfloor$ and $r=n-2^t$.  For $n=1$ use $(1)$.  For
$n>1$, define
\begin{equation}\label{eq:witness}
 W_n=
 \begin{cases}
 C_t,&r=0,\\
 (r,C_t),&r=2^{t-1},\\
 (C_t,r),&0<r<2^t,\ r\ne2^{t-1}.
 \end{cases}
\end{equation}
If $0<r<2^t$ and $r\ne2^{t-1}$, then for $0\le j<t$ the orbit is
\[
 A_n^j(W_n)=(2^j,2^j,2^{j+1},\ldots,2^{t-1},r),
\]
where the power string after the repeated pair is empty when $j=t-1$.
Indeed, the appended $r$ never equals its unchanged left neighbour
$2^{t-1}$.  The $t$th state is $(2^t,r)$, which is fixed.

If $r=2^{t-1}$ and $t\ge2$, then for $0\le j<t-1$,
\[
 A_n^j(W_n)=(r,2^j,2^j,2^{j+1},\ldots,2^{t-1}),
\]
and $A_n^{t-1}(W_n)=(r,r,r)$, followed by the fixed state $(3r)$.
The remaining bases are explicit: $W_1=(1)$ has depth zero,
$W_2=C_1=(1,1)$ has depth one, and the half-remainder case
$W_3=(1,1,1)$ also has depth one.  Thus $W_n$ has depth $t$ in every case.
This proves \eqref{eq:clock} and the equality statement.

\section{The divisor-path inverse}

\begin{proof}[Proof of Theorem~\ref{thm:main}(3)]
Let $\alpha$ map to $\beta=(b_1,\ldots,b_k)$.  The unique maximal run of
$\alpha$ producing $b_i$ has some common value $s_i$ and length $b_i/s_i$;
therefore $s_i\mid b_i$.  Consecutive runs are maximal, so
$s_i\ne s_{i+1}$.  The source length is $\sum_i b_i/s_i$.

Conversely, choose a divisor $s_i\mid b_i$ for every $i$, with adjacent
choices unequal, and expand $b_i$ into $b_i/s_i$ copies of $s_i$.  The
adjacent inequality makes these expanded blocks precisely the maximal runs,
so the expansion maps to $\beta$.  The two constructions are inverse.  Sorting
the choices by total expanded length proves \eqref{eq:coefficient} and all its
consequences.
\end{proof}

Formula \eqref{eq:fibrepoly} is also an efficient exact atlas.  At layer $i$
one keeps the last divisor $s_i$ and a polynomial in the accumulated source
length; transitions are allowed between unequal divisors.  Thus every target
fibre can be computed by a weighted path dynamic program over its divisor
sets, without enumerating all $2^{n-1}$ source compositions.

For completeness, if $C(x)$ includes the empty fixed composition, the
classical fixed-class generating function follows from a one-line last-part
decomposition.  If $C_j(x)$ counts nonempty fixed compositions ending in
$j$, then $C_j=x^j(C-C_j)$, whence
\begin{equation}\label{eq:carlitz}
 C(x)=\frac{1}{1-\displaystyle\sum_{j\ge1}\frac{x^j}{1+x^j}}.
\end{equation}
Equation \eqref{eq:carlitz} is included only to identify the recurrent class;
it receives no contribution credit.

\section{Exact falsification and boundary}

The accompanying standard-library program enumerates every positive
composition of every total $1\le n\le18$.  In a fresh process it checks total
preservation, strict descent, absence of nontrivial cycles, every endpoint,
the pointwise logarithmic bound, the witness \eqref{eq:witness}, the classical
fixed census, and \eqref{eq:fibrepoly} for every possible target, including
empty fibres.  The computation covers $262{,}143$ states and executes
$2{,}690{,}869$ exact assertions.  For $n=18$ the depth profile is
\[
 10843+60946z+47626z^2+11415z^3+242z^4.
\]
These finite checks are counterexample pressure and reproducibility evidence,
not proofs of the all-parameter statements.

The closest occupied internal carrier is a composition-refinement system, but
that map splits balanced blocks and uses a suffix decoder.  Here the operation
coarsens maximal equal runs, the clock is forced by weight-doubling ancestry,
and the inverse is arithmetic on adjacent divisor sets.  Ordinary run-length
encoding likewise records a value/count pair and changes representation;
$A_n$ multiplies the pair, preserves total weight, and iterates inside the same
finite carrier.  Recent work on random weak-composition growth studies the
evolution and disappearance of adjacent equalities~\cite{BevanThrelfall2025},
while recent Arndt--Carlitz work remains a static adjacent-restriction
study~\cite{HopkinsTangboonduangjit2026}; both receive zero contribution
credit and use different objects.  A bounded source search found no direct
owner for this exact conjunction, but a search non-hit is not a novelty or
priority certificate.

\section*{Data availability}
No external data were used.  The exact verifier, frozen transcript, and claim
ledger are included with this manuscript.

\section*{Ethics declaration}
This mathematical study involved no human participants, animals, personal
data, or field intervention; ethics approval was not applicable.

\section*{Author contributions}
For anonymous internal review, a single anonymous author is recorded as
responsible for conceptualization, formal analysis, software, validation,
writing, and reproducibility curation.

\section*{Conflict of interest}
No conflict of interest is declared.

\section*{Funding}
No external funding is declared for this anonymous internal note.

\section*{Limitations and external status}
The owner search was bounded and terminology dependent.  The paper does not
claim novelty for Carlitz compositions or run statistics, and it does not
analyze random initial states, asymptotic depth distributions, or fibres of
higher iterates.  All public posting, submission, priority, authorship, and
release actions remain under \texttt{HOLD\_EXTERNAL}.

\bibliographystyle{plainnat}
\bibliography{references}

\end{document}
